\section{Erd\H{o}s Problem 182: the number of edges forcing a regular subgraph}

\subsection{1) Statement and scope}
Fix $k\geq3$, and let $E_k(n)$ be the largest number of edges in an
$n$-vertex graph with no nonempty $k$-regular subgraph. The subgraph
need not be induced or spanning. Is $E_k(n)=n^{1+o(1)}$?
The known answer is the more precise order
\[
 E_k(n)=\Theta_k(n\log\log n).
\]
We deduce this from two explicitly stated external results and give a
self-contained common-neighbour bound explaining an elementary baseline.

\subsection{2) The matching external results}
Janzer--Sudakov prove that, for each fixed $k$, some $C_k$ has the
following property: every sufficiently large $n$-vertex graph of average
degree at least $C_k\log\log n$ contains a nonempty $k$-regular subgraph.
Changing the base of the logarithms only changes the constants and a
bounded lower-order term. The handshake identity therefore gives
\[
 E_k(n)<\tfrac12 C_kn\log\log n.                       \tag{1}
\]
For the reverse direction, the construction of Pyber--R\"odl--Szemer\'edi,
stated as Theorem 1.1 in the same paper, supplies for all sufficiently
large $n$ a graph with at least $cn\log\log n$ edges and no regular
subgraph of any degree at least three. In particular it has no
$k$-regular subgraph. Thus
\[
 E_k(n)\geq cn\log\log n.                              \tag{2}
\]
Both the construction and the regular-subgraph theorem are external
inputs; the assertion for every degree at least three in the lower
construction is essential and is part of its stated conclusion.

Equations (1)--(2) prove the claimed order. For each $\epsilon>0$,
$\log\log n=o(n^\epsilon)$, and also
$\log(n\log\log n)/\log n\to1$. This verifies both the requested
upper bound $n^{1+o(1)}$ and the exponent-one interpretation of the
answer, with $k$ fixed before $n$ tends to infinity.

\subsection{3) A direct, weaker bound from common neighbours}
A $K_{k,k}$ is itself $k$-regular, so a graph counted by $E_k(n)$ must
be $K_{k,k}$-free. Every $k$-set of vertices has at most $k-1$ common
neighbours, all outside that set. Double counting gives
\[
 \sum_{v}\binom{d(v)}k\leq(k-1)\binom nk.
\]
For an integer $d\geq0$,
$\binom dk\geq(d-k+1)_+^k/k!$, where $u_+=\max(u,0)$.
The function $u\mapsto(u-k+1)_+^k$ is convex. Writing
$\bar d=2e/n$ and applying Jensen's inequality, we obtain
\[
 \frac n{k!}(\bar d-k+1)_+^k
 \leq(k-1)\frac{n^k}{k!}.
\]
Hence
\[
 e\leq\frac12(k-1)n+\frac12(k-1)^{1/k}n^{2-1/k}.
\]
This argument is complete, but its exponent is too large to replace
the external theorem in Section 2. It identifies precisely the loss
in restricting attention to complete bipartite regular subgraphs.

\subsection{4) Outcome and sources}
\textbf{Outcome: SOLVED using the stated established upper and lower
theorems.} The degree conversion, asymptotic quantifiers, and elementary
comparison are supplied here. Neither deep theorem is independently
reproved, and no new bound is claimed.
\begin{itemize}
\item O. Janzer and B. Sudakov, \emph{Resolution of the Erd\H{o}s--Sauer
problem on regular subgraphs}, Theorems 1.1 and 1.2,
\texttt{https://arxiv.org/abs/2204.12455}.
\item T. F. Bloom, Problem 182,
\texttt{https://www.erdosproblems.com/182}, checked 23 September 2026.
\end{itemize}
The percentage includes the two explicitly identified external inputs.
\subsection{5) COMPLETION ESTIMATE}
\textbf{COMPLETION: 100\%}
